% from ega4-16-fr.tex
\oldpage[IV-4]{45}
comme $B$-module à droite par son second facteur $B$ et
$(B\otimes_A B)/\mathfrak{I}^{n+1}$ comme $B$-module à gauche par son premier
facteur $B$ (16.7.2). Il revient au même de définir un homomorphisme de
$B$-modules
\[
  \varphi_0:B\otimes_A B\longrightarrow
  ((B\otimes_A B)/\mathfrak{I}^{n'+1})
  \otimes_B((B\otimes_A B)/\mathfrak{I}^{n+1})
\]
et de prouver qu'il s'annule dans $\mathfrak{I}^{n+n'+1}$. Or, on définit
aussitôt un tel homomorphisme par la condition que
\[
  \varphi_0(b\otimes b')
  =\pi_{n'}(b\otimes1)\otimes\pi_n(1\otimes b')
  \qquad\text{pour $b,b'$ dans $B$}
\]
avec les notations de (16.3.7). En outre, il est immédiat que $\varphi_0$ est
un homomorphisme d'\emph{anneaux}. Or, on peut écrire
\[
  \begin{aligned}
  \varphi_0(b\otimes1-1\otimes b)
  ={}&\pi_{n'}(b\otimes1-1\otimes b)\otimes\pi_n(1\otimes1)\\
  &+\pi_{n'}(1\otimes b)\otimes\pi_n(1\otimes1)
  -\pi_{n'}(1\otimes1)\otimes\pi_n(1\otimes b)
  \end{aligned}
\]
et on a
\[
  \begin{aligned}
  \pi_{n'}(1\otimes b)\otimes\pi_n(1\otimes1)
  &=\pi_{n'}(1\otimes1)b\otimes\pi_n(1\otimes1)\\
  &=\pi_{n'}(1\otimes1)\otimes b\pi_n(1\otimes1)\\
  &=\pi_{n'}(1\otimes1)\otimes\pi_n(b\otimes1)
  \end{aligned}
\]
d'où finalement
\[
  \begin{aligned}
  \varphi_0(b\otimes1-1\otimes b)
  ={}&\pi_{n'}(b\otimes1-1\otimes b)\otimes\pi_n(1\otimes1)\\
  &+\pi_{n'}(1\otimes1)\otimes\pi_n(b\otimes1-1\otimes b).
  \end{aligned}
  \tag{16.8.9.4}
  \label{IV.16.8.9.4}
\]
Un produit de $n+n'+1$ termes de la forme (16.8.9.4) est donc
nécessairement nul, car il en est ainsi d'un produit de $n+1$ termes de la
forme $\pi_n(b\otimes1-1\otimes b)$ et d'un produit de $n'+1$ termes de la
forme $\pi_{n'}(b\otimes1-1\otimes b)$. La conclusion résulte donc de
(0, 20.4.4).

\begin{corollary}[16.8.10]
\label{IV.16.8.10}
Le faisceau $\mathscr{D}\!iff_{X/S}(\mathcal{O}_X,\mathcal{O}_X)$ (aussi noté
$\mathscr{D}\!iff_{X/S}$) est canoniquement muni d'une structure de faisceau
d'anneaux, les $\mathscr{D}\!iff_{X/S}^n$ formant une filtration croissante
compatible avec cette structure.
\end{corollary}

En particulier, $\mathscr{D}\!iff_{X/S}^0$ est un faisceau de sous-anneaux de
$\mathscr{D}\!iff_{X/S}$, qui s'identifie canoniquement à $\mathcal{O}_X$
(16.8.1). Les formules (16.8.5.1) et (16.8.5.2) montrent que la structure de
$\mathcal{O}_X$-Bimodule de $\mathscr{D}\!iff_{X/S}$ provient de la
multiplication à gauche et à droite par les sections de $\mathcal{O}_X$
considéré comme faisceau de sous-anneaux de $\mathscr{D}\!iff_{X/S}$.

\begin{remarks}[16.8.11]
\label{IV.16.8.11}
\begin{enumerate}
  \item[(i)] Supposons que
  $\mathscr{F}=\bigoplus_{\lambda\in L}\mathscr{F}_\lambda$; alors il est
  clair (16.7.2.1) que
  $\mathscr{P}_{X/S}^n(\mathscr{F})=
  \bigoplus_{\lambda\in L}\mathscr{P}_{X/S}^n(\mathscr{F}_\lambda)$; comme le
  foncteur $\mathscr{F}\mapsto\Gamma(U,\mathscr{F})$ commute à la formation
  des sommes directes quelconques, $d_{X/S,\mathscr{F}}^n$ est
  l'homomorphisme dont la restriction à chaque $\mathscr{F}_\lambda$ est
  \[
    d_{X/S,\mathscr{F}_\lambda}^n:\mathscr{F}_\lambda
    \longrightarrow\mathscr{P}_{X/S}^n(\mathscr{F}_\lambda);
  \]
  on en conclut aussitôt que l'on a
  \[
    \operatorname{Diff}_{X/S}^n(\mathscr{F},\mathscr{G})
    =\prod_{\lambda\in L}
    \operatorname{Diff}_{X/S}^n(\mathscr{F}_\lambda,\mathscr{G}),
  \]
  et par suite aussi (0\textsubscript{I}, 3.2.6)
  \[
    \mathscr{D}\!iff_{X/S}^n(\mathscr{F},\mathscr{G})
    =\prod_{\lambda\in L}
    \mathscr{D}\!iff_{X/S}^n(\mathscr{F}_\lambda,\mathscr{G}).
  \]
  Par ailleurs, si $\mathscr{G}=\prod_{\mu\in M}\mathscr{G}_\mu$
  (0\textsubscript{I}, 3.2.6), on a
  \[
    \operatorname{Hom}_{\mathcal{O}_X}
    (\mathscr{P}_{X/S}^n(\mathscr{F}),\mathscr{G})
    =\prod_{\mu\in M}\operatorname{Hom}_{\mathcal{O}_X}
    (\mathscr{P}_{X/S}^n(\mathscr{F}),\mathscr{G}_\mu),
  \]
  